\section{Fractional Knapsack Problem}
\label{sec:fractional-knapsack}

\problemstatement{We are given $n$ items, where the $i$-th item has a value
$\textit{val}[i]$ and a weight $\textit{wt}[i]$, and a knapsack of capacity
$W$. Unlike the classical $0/1$ knapsack, here we are allowed to take any
\emph{fraction} of an item -- if we take a fraction $x \in [0,1]$ of item $i$,
we gain $x \cdot \textit{val}[i]$ value while consuming $x \cdot
\textit{wt}[i]$ weight. We must choose fractions for each item, subject to
the total weight taken not exceeding $W$, so as to \textbf{maximize the total
value} carried out of the knapsack.}

\begin{patternbox}
The single word \emph{``fraction''} (or \emph{``can be broken into
smaller units''}) attached to a knapsack-style capacity problem is the
loudest signal in this entire chapter: it rules out dynamic programming
immediately. The $0/1$ knapsack needs DP precisely because an item is
all-or-nothing, which creates genuine overlapping subproblems (``did we take
item $i$ or not''). The moment items become divisible, that branching
disappears -- we can always characterize the best item to put in next using a
single number, its \emph{value-per-unit-weight ratio}, and greedily fill the
knapsack from the best ratio downward. Any time you see ``maximize total
value/profit under a single weight/capacity/budget constraint, items may be
split,'' reach for the ratio-greedy pattern.
\end{patternbox}

\subsection*{Understanding the problem carefully}

Define, for each item $i$, its \emph{rate of return}
\[
  r_i = \frac{\textit{val}[i]}{\textit{wt}[i]}.
\]
This number tells us exactly how much value we earn per unit of weight we
spend if we take item $i$. Because we are allowed to take fractions, taking
item $i$ is like having access to an unlimited (well, limited by
$\textit{wt}[i]$) supply of ``value at rate $r_i$.'' If our only constraint is
a total weight budget $W$, then to maximize value, common sense says: spend
the budget first on the item with the best rate of return, then the next
best, and so on, until the budget runs out -- possibly using only part of the
last item we touch.

\begin{ideabox}[Greedy Choice]
Sort the items in decreasing order of $r_i = \textit{val}[i] /
\textit{wt}[i]$. Process them in this order; for each item, take as much of
it as the remaining capacity allows -- the whole item if it fits, otherwise
exactly the fraction of it that exactly fills the remaining capacity, after
which the knapsack is full and we stop.
\end{ideabox}

\subsection*{Why the greedy choice is correct}

Suppose, for contradiction, that some optimal solution does not follow this
rule -- say it takes a positive amount $a > 0$ of a lower-ratio item $i$
while leaving some amount $b > 0$ of a higher-ratio item $i'$ (with $r_{i'} >
r_i$) untaken, even though there was room for it. Then we can perform a small
exchange: reduce the amount of item $i$ taken by $\varepsilon =
\min(a, b)$ in weight, and increase the amount of item $i'$ taken by the same
weight $\varepsilon$. The total weight used is unchanged, so the solution
remains feasible. But the value changes by
\[
  \Delta = \varepsilon \cdot r_{i'} - \varepsilon \cdot r_i
         = \varepsilon (r_{i'} - r_i) > 0,
\]
since $r_{i'} > r_i$. So the exchange strictly increases total value, which
contradicts the assumption that the original solution was optimal. Hence no
optimal solution can avoid fully preferring higher-ratio items over
lower-ratio ones whenever weight can be moved between them -- which is
exactly the greedy rule. Because fractional amounts are allowed, this
exchange can always be performed down to an infinitesimal degree, so the
greedy strategy is not just \emph{a} good heuristic here, it is provably
optimal, with no gap to close via DP.

\subsection*{Algorithm}

\begin{enumerate}
  \item Compute the ratio $r_i = \textit{val}[i] / \textit{wt}[i]$ for every
        item.
  \item Sort the items in decreasing order of $r_i$.
  \item Initialize $\textit{remaining} \gets W$ and
        $\textit{totalValue} \gets 0$.
  \item For each item $i$ in sorted order:
  \begin{enumerate}
    \item If $\textit{wt}[i] \le \textit{remaining}$: take the whole item.
          $\textit{totalValue} \gets \textit{totalValue} + \textit{val}[i]$;
          $\textit{remaining} \gets \textit{remaining} - \textit{wt}[i]$.
    \item Otherwise: take the fraction
          $f = \textit{remaining} / \textit{wt}[i]$ of this item.
          $\textit{totalValue} \gets \textit{totalValue} + f \cdot
          \textit{val}[i]$; set $\textit{remaining} \gets 0$ and stop --
          the knapsack is now full.
  \end{enumerate}
  \item Return $\textit{totalValue}$.
\end{enumerate}

\begin{algorithm}[H]
\caption{Fractional Knapsack}
\begin{algorithmic}[1]
\Procedure{FractionalKnapsack}{$\textit{val}, \textit{wt}, W$}
  \State sort items by $\textit{val}[i]/\textit{wt}[i]$ in decreasing order
  \State $\textit{remaining} \gets W$, $\textit{totalValue} \gets 0$
  \For{each item $i$ in sorted order}
    \If{$\textit{remaining} = 0$}
      \State \textbf{break}
    \EndIf
    \If{$\textit{wt}[i] \le \textit{remaining}$}
      \State $\textit{totalValue} \gets \textit{totalValue} + \textit{val}[i]$
      \State $\textit{remaining} \gets \textit{remaining} - \textit{wt}[i]$
    \Else
      \State $f \gets \textit{remaining} / \textit{wt}[i]$
      \State $\textit{totalValue} \gets \textit{totalValue} + f \times \textit{val}[i]$
      \State $\textit{remaining} \gets 0$
    \EndIf
  \EndFor
  \State \Return \textit{totalValue}
\EndProcedure
\end{algorithmic}
\end{algorithm}

\subsection*{Dry run}

\dryrun{Take three items with $(\textit{val}, \textit{wt})$ pairs
$(60, 10)$, $(100, 20)$, $(120, 30)$, and knapsack capacity $W = 50$.}

First compute ratios:
\[
  r_1 = 60/10 = 6.0, \qquad r_2 = 100/20 = 5.0, \qquad r_3 = 120/30 = 4.0.
\]
Sorted in decreasing order of ratio, the items are already in the order
$1, 2, 3$.

\begin{itemize}
  \item \textbf{Item 1} ($\textit{val}=60, \textit{wt}=10$): remaining
        capacity is $50$, and $10 \le 50$, so we take all of it.
        $\textit{totalValue} = 60$, $\textit{remaining} = 40$.
  \item \textbf{Item 2} ($\textit{val}=100, \textit{wt}=20$): remaining
        capacity is $40$, and $20 \le 40$, so we take all of it.
        $\textit{totalValue} = 160$, $\textit{remaining} = 20$.
  \item \textbf{Item 3} ($\textit{val}=120, \textit{wt}=30$): remaining
        capacity is $20$, but $30 > 20$, so we cannot take the whole item.
        We take the fraction $f = 20/30 = 2/3$. This contributes
        $\frac{2}{3} \times 120 = 80$ value.
        $\textit{totalValue} = 160 + 80 = 240$, $\textit{remaining} = 0$.
\end{itemize}

The knapsack is now full, so we stop. The maximum value obtainable is
$\boxed{240}$.

\subsection*{C++ implementation}

\begin{lstlisting}
#include <bits/stdc++.h>
using namespace std;

struct Item {
    int value;
    int weight;
};

double fractionalKnapsack(int W, vector<Item>& items) {
    // Sort items by value/weight ratio in decreasing order
    sort(items.begin(), items.end(), [](const Item& a, const Item& b) {
        double r1 = (double)a.value / (double)a.weight;
        double r2 = (double)b.value / (double)b.weight;
        return r1 > r2;
    });

    int remaining = W;
    double totalValue = 0.0;

    for (auto& item : items) {
        if (remaining == 0) break;

        if (item.weight <= remaining) {
            totalValue += item.value;
            remaining -= item.weight;
        } else {
            double fraction = (double)remaining / (double)item.weight;
            totalValue += fraction * item.value;
            remaining = 0;
        }
    }
    return totalValue;
}

int main() {
    vector<Item> items = {{60, 10}, {100, 20}, {120, 30}};
    int W = 50;

    cout << "Maximum value: " << fractionalKnapsack(W, items) << endl;
    return 0;
}
\end{lstlisting}

\begin{complexitybox}
\textbf{Time Complexity.} Sorting the $n$ items by ratio costs
$O(n \log n)$. The subsequent single pass over the sorted items costs
$O(n)$. Hence the total time complexity is
\[
  O(n \log n).
\]
\textbf{Space Complexity.} Sorting is performed in place (or with a small
$O(n)$ auxiliary array depending on implementation), and we use only a
constant number of extra scalar variables, giving auxiliary space
complexity $O(1)$ beyond the input storage (or $O(n)$ if we count the
input array itself, and $O(\log n)$ for the sort's recursion stack).
\end{complexitybox}

\begin{takeawaybox}
Divisibility is the tell. The instant a knapsack-style problem allows
\emph{fractions} of an item, the all-or-nothing branching that forces
dynamic programming in $0/1$ knapsack disappears, and a single
value-per-weight ratio becomes a perfect, provably optimal greedy ranking.
Contrast this sharply with Section~\ref{sec:assign-cookies}: there, sorting
gave us a feasibility-matching greedy; here, sorting by ratio gives us a
value-maximizing greedy, but the underlying exchange-argument justification
is the same family of reasoning.
\end{takeawaybox}
